\section{A framework based on finite presentations and proof transport}
\label{sec:framework}

\subsection{The arithmetic base and the meaning of an embedding}

Our base theory is first-order intuitionistic logic with equality,
one function symbol for each primitive recursive definition, their
defining equations, successor injectivity and nonzeroness, and induction
for quantifier-free formulas, with parameters. This is the presentation
implemented by the repository's \code{PR}, \code{PTm}, \code{PRAAx}
and \code{Prf} types \cite{PRSource,DeductionSource}. We do not add
induction for an arbitrary universally quantified assertion, excluded
middle for arbitrary propositions, choice, or an axiom of convergence.

The term ``embedding'' here means a specified translation of the finite
content of a mathematical statement into arithmetic. It does not mean
that every theorem of a stronger ambient theory becomes a theorem of
\ipra. The translation records the representation, the observations
being requested, a concrete algorithm, and a derivation of its stated
finite contract. Distinct translations can preserve different information.
A bound on the number of points, a list of all points, a rational
approximation and a coherent family of approximations are distinct outputs.

Use any explicitly fixed PR pairing and list coding. For definiteness
one may pair $a,b$ by $(a+b)(a+b+1)/2+b$, and encode a nonempty list
by one plus the pair of its head and tail. Projections are bounded
searches; a nonzero tail code decreases. All list folds therefore have
PR definitions with fuel at most the input code. More efficient list
implementations are extensionally interchangeable after their simulation
law has been proved. Encode integers by a sign bit and magnitude, and
rationals by a signed numerator and a positive denominator, normalized
by gcd; zero has denominator one. A malformed code has a designated
default and a separate validity flag. A theorem never relies on the
default satisfying the specification of a valid mathematical object.

Rational equality, order and arithmetic reduce to signed integer
arithmetic and positive-denominator cross multiplication. The Euclidean
algorithm can be run for $a+b+1$ steps on nonnegative inputs $a,b$,
freezing the state when the second component becomes zero. Each active
step strictly decreases a nonnegative component. Its invariant tracks
the common divisors and, for extended gcd, two linear-combination
identities. Thus division, gcd, lcm and rational normalization do not
require unbounded minimization. Sparse multivariate polynomials are
finite lists of coefficient/exponent-vector pairs; evaluation is a fixed
fold using finite powers. These concrete data operations will also be
defined in the Lean target file.

\subsection{Bounded formulas and allowed induction}

For a fixed PR Boolean predicate $P(a,i)$, put
\[
 A(a,0)=1,\qquad A(a,t+1)=A(a,t)\mathbin{\wedge}P(a,t).
\]
Then a bounded assertion $\forall i<t\,P(a,i)=1$ is represented by
the atomic equality $A(a,t)=1$. A corresponding disjunction fold
represents bounded existence. The familiar unfolding equivalences
are proved with the queried index free. For the converse under a fixed
bound $B$, use the guarded quantifier-free invariant
$t\le B\to A(a,t)=1$, with the bounded-universal premise kept
outside the induction formula. For nested bounded
quantifiers the construction is repeated, from the inside outward.
This is a schema for each fixed formula and its PR predicates.

Finite-list correctness assertions, such as ``each processed coefficient
agrees'' or ``all entries satisfy this divisibility equation'', can
therefore be used as quantifier-free invariants. When a proof concerns
an arbitrary solution $p$, keep $p$ as a free parameter during induction.
Do not put ``for every rational point'' inside the induction invariant.
This simple distinction matters in the enumeration proof: the finite
matrix calculation is performed with one arbitrary point fixed, and
universal generalization occurs afterwards.

There is no total PR interpreter for arbitrary PR program codes. If
one had $E(e,x)$ and every unary PR function were represented, the
PR function $d(x)=E(x,x)+1$ would have a code $e$ and give
$d(e)=d(e)+1$. Our constructions avoid that contradiction. The simulation
theorems are schemas for specified PR syntax. The polynomial checker
is uniform over arbitrary polynomial \emph{data}, because it interprets
one fixed polynomial grammar with one fixed PR evaluation algorithm.
The two forms of uniformity are compatible and must be stated separately.

\subsection{Finite presentations and observations}

A finite presentation has a validity predicate $V(a)$ and, when needed,
a decidable equivalence relation $a\sim b$ on valid codes. A restricted
operation has a total PR implementation and a separate guard under which
its mathematical law holds. Examples include division with a nonzero
denominator, a matrix inverse with a verified inverse table, and a
localization coordinate with a nonzero polynomial value.

An observation is a fixed PR function or relation on the representation.
It may be an exact polynomial coefficient, a matrix entry, a residue
modulo $q^n$, a rational number with an error bound, or membership in a
finite combinatorial object. A mathematical proof is organized by the
observations it actually uses. Infinite terminology is explanatory until
the corresponding arithmetic observations and their precision demands
have been specified.

For example, a metric name $a(k)\in\Q$ can be required to satisfy
$|a(k)-a(\ell)|\le2^{-k}+2^{-\ell}$ at the finitely many indices
needed in a calculation. A formal series is observed by exact coefficient
truncation. A residue tower is observed modulo a specified power. These
are different contracts. There is no rule converting a formal series
into an analytic function without a tail estimate, or a local residue
tower into a rational point without a global argument.

\subsection{Syntactic iteration fusion}

For a specified unary PR operation $f$, define
\[
 \iter_f(0,x)=x,\qquad
 \iter_f(n+1,x)=f(\iter_f(n,x)).
\]
This definition is itself primitive recursion. In the repository syntax,
the recursion counter comes first and the previous result is the second
input to the step. The Lean construction chooses exactly those projections.

\begin{theorem}[Proof-preserving fusion]\label{thm:fusion}
Let $f,g,h$ be specified unary PR function symbols. If an actual
\ipra\ derivation of $h(f(x))=g(h(x))$ is supplied, a finite
syntactic construction produces an actual \ipra\ derivation of
\[
 \forall n\;h(\iter_f(n,x))=\iter_g(n,h(x)).
 \tag{F}\label{eq:fusion}
\]
\end{theorem}

\begin{proof}
Use the equality in \eqref{eq:fusion} as the quantifier-free invariant
$I(n,x)$. At zero, the two recursion equations give $h(x)=h(x)$.
For the successor step, start with the induction hypothesis and compute
inside the object equality calculus:
\begin{align*}
 h(\iter_f(n+1,x))
 &=h(f(\iter_f(n,x)))\\
 &=g(h(\iter_f(n,x)))\\
 &=g(\iter_g(n,h(x)))\\
 &=\iter_g(n+1,h(x)).
\end{align*}
The second equality is the supplied local derivation after substituting
$\iter_f(n,x)$ for its variable; the third is equality congruence
applied to the induction hypothesis. The first and last use defining
equations. Apply the object rule for quantifier-free induction, keeping
$x$ free. The operations used on derivations are substitution,
weakening, congruence, transitivity and the induction constructor.
Every one is an existing operation on \code{Prf}; no semantic reflection
step occurs. \code{ObjectFusion.fusion} implements this proof.
\end{proof}

The theorem also covers algorithms with numerical parameters. Put the
parameter in the state and use the unary step $(a,x)\mapsto(a,F(a,x))$.
A representation map can preserve that parameter. This is a fixed
PR pair operation, not a family of programs selected by a universal
interpreter. A step that depends on its iteration index is handled by
also storing and incrementing the counter in the state.

\subsection{Relations, guards and simulation}

Equality is often too restrictive. Quotient representatives may differ,
and rational approximations carry inequalities instead of exact equalities.
The relational theorem gives the appropriate common statement.

\begin{theorem}[Quantifier-free simulation]\label{thm:simulation}
Let $R(x,y)$ be a specified quantifier-free arithmetic formula and let
$f,g$ be specified unary PR steps. From a derivation
\[
                      R(x,y)\ \vdash\ R(f(x),g(y))
\]
one can construct a derivation
\[
       R(x,y)\ \vdash\ \forall n\,
                      R(\iter_f(n,x),\iter_g(n,y)).
 \tag{S}\label{eq:simulation}
\]
\end{theorem}

\begin{proof}
Substitution into a quantifier-free formula preserves its quantifier-free
status, by structural inspection of the formula. Hence
$R(\iter_f(n,x),\iter_g(n,y))$ is an allowed induction invariant.
At zero, the hypothesis $R(x,y)$ and the two defining equations give
the base case. Substitute the two current iterates for $x,y$ in the
supplied local derivation. Its premise is the induction hypothesis;
its conclusion becomes the next invariant after applying the successor
equations. Quantifier-free induction proves \eqref{eq:simulation}.
The formal proof keeps the shifted initial premise in the context under
the counter binder. \code{ObjectSimulation.simulation} checks these
substitution and variable-scope details explicitly.
\end{proof}

Guards should be part of $R$. For example,
$R(x,y)$ may assert that both codes are valid and that their evaluations
agree, or that both denominators are positive and an integer-cleared
error inequality holds. A local lemma valid only on well-formed inputs
then preserves the entire invariant. This prevents a later iteration
from invoking a law outside its domain.

\begin{corollary}[Transport through repeated squaring]
\label{cor:squaring}
Suppose a specified map $h$ between two coded multiplication systems
satisfies the object derivation
$h(a\cdot b)=h(a)\star h(b)$.
Define $s_0=x$, $s_{n+1}=s_n\cdot s_n$, and
$t_0=h(x)$, $t_{n+1}=t_n\star t_n$. Then \ipra\ proves
$h(s_n)=t_n$ for every $n$.
\end{corollary}
\begin{proof}
Substitute $a=b=x$ into the local multiplication derivation to obtain
the premise of Theorem~\ref{thm:fusion} for the two squaring steps.
Apply that theorem. No associativity is needed for this particular
statement; an identification with $x^{2^n}$ may be proved separately
in an associative system. The supplied Lean corollary performs the
substitution and invokes the checked fusion builder.
\end{proof}

The same idea applies to a polynomial expression evaluated by an
instruction list: store its program counter, stack and fixed parameters
in the state. A local stack-operation simulation yields correctness of
the whole finite evaluator. This is an interpretation of a fixed finite
expression grammar, whose operations are specified; it is not an
interpreter for all PR functions.

\subsection{Exact witnesses from a finite observable image}

The algebraic theorem in the next section supplies a finite grid
containing all possible rational values of selected observables. Here
is the general step that turns such a grid into object witnesses.

\begin{proposition}[Finite reconstruction]\label{prop:reconstruction}
Let $V(p)$ be a decidable predicate on finite codes, let
$o(p)$ be a fixed PR observation, and let $G$ be an explicitly computed
finite list. Suppose the following statements have \ipra\ proofs:
\begin{enumerate}
\item $V(p)$ implies $o(p)\in G$;
\item for each $q\in G$, a fixed PR list producer $B(q)$ contains
every valid $p$ with $o(p)=q$;
\item membership in $B(q)$ and all validity and observation tests are
finite decidable tests.
\end{enumerate}
Then filter $\bigcup_{q\in G}B(q)$ by $V$ and remove duplicates.
The resulting PR list contains exactly the valid objects.
\end{proposition}
\begin{proof}
A retained object satisfies $V$ by its filter test. Conversely, fix an
arbitrary valid $p$. The first premise supplies its observation among
the finite $G$; the second supplies $p$ in that observation's list.
Therefore $p$ is visited and retained. A finite deduplication pass
preserves membership. Its invariant equates membership in the output
prefix with membership in the processed input prefix; bounded folds
make that invariant quantifier-free. This proves both directions.
\end{proof}

When the observation is the full rational coordinate tuple, $B(q)$
is the singleton $[q]$. This is the important algebraic specialization.
When an observation is constant on an infinite family, there is no
finite reconstruction list in general. The framework exposes the
missing premise instead of concealing it in the word ``representation''.
